% !TeX root = ../main.tex
%%%%%%%%%%%%%%%%%%%%%%%%%%   3   %%%%%%%%%%%%%%%%%%%%%%%%%%%%%%
\section{后期拟完成的研究工作及进度安排}

后续研究拟在现有理论整理与基线实验的基础上，围绕神经网络与非线性滤波的结合，分两个阶段推进：

\begin{enumerate}[label=(\arabic*),leftmargin=2.5em,topsep=3pt,itemsep=6pt,parsep=0pt]
  \item \textbf{2026 年 9--10 月：文献学习与理论总结。}
  学习神经网络相关的滤波论文，重点梳理网络在状态预测、后验更新及输运映射中的作用，分析模型假设、优化目标与算法步骤。总结代表性方法的数学原理、适用条件和局限性，形成文献笔记与算法对照，为后续设计确定研究切入点。
  \item \textbf{2026 年 11--12 月：仿真环境选择与算法验证。}
  根据前期分析筛选合适的仿真环境，明确状态转移、观测模型、噪声设置与评价指标。针对选定问题设计并实现相应的滤波算法，开展必要的训练、基线对比与消融实验，分析估计精度、计算成本和稳定性，形成“问题建模—算法设计—仿真验证—结果分析与改进”的研究闭环。
\end{enumerate}
